% !TeX encoding = UTF-8
\documentclass[14pt,a4paper]{extarticle}

\newif\ifusekprj
\IfFileExists{kprj.sty}{\usekprjtrue\usepackage{kprj}}{\usekprjfalse}
\usepackage{tabularx}
\usepackage{float}
\usepackage{flafter}
\usepackage{graphicx}
\usepackage{booktabs}
\usepackage{longtable}
\usepackage{multirow}
\usepackage{array}
\usepackage{amsmath,amssymb}
\usepackage[T2A]{fontenc}
\usepackage[utf8]{inputenc}
\usepackage[russian]{babel}
\usepackage{hyperref}
\usepackage{listings}
\usepackage{xcolor}

\makeatletter
\ifusekprj
\else
  \newcommand{\authorfull}[1]{\gdef\@authorfull{#1}}
  \newcommand{\faculty}[1]{\gdef\@faculty{#1}}
  \newcommand{\facultyshort}[1]{\gdef\@facultyshort{#1}}
  \newcommand{\chair}[1]{\gdef\@chair{#1}}
  \newcommand{\chairno}[1]{\gdef\@chairno{#1}}
  \newcommand{\group}[1]{\gdef\@group{#1}}
  \newcommand{\chief}[1]{\gdef\@chief{#1}}
  \newcommand{\FirstConsultant}[1]{\gdef\@FirstConsultant{#1}}
  \newcommand{\inspector}[1]{\gdef\@inspector{#1}}
  \newcommand{\workyear}[1]{\gdef\@workyear{#1}}
  \newcommand{\setcontentsname}[1]{\renewcommand{\contentsname}{#1}}
  \newcommand{\setrefsname}[1]{\gdef\@refsname{#1}}
  \providecommand{\@authorfull}{}
  \providecommand{\@faculty}{}
  \providecommand{\@facultyshort}{}
  \providecommand{\@chair}{}
  \providecommand{\@chairno}{}
  \providecommand{\@group}{}
  \providecommand{\@chief}{}
  \providecommand{\@FirstConsultant}{}
  \providecommand{\@inspector}{}
  \providecommand{\@workyear}{\the\year}
  \providecommand{\@refsname}{СПИСОК ИСПОЛЬЗОВАННЫХ ИСТОЧНИКОВ}
\fi
\makeatother

\hypersetup{
  colorlinks=true,
  linkcolor=black,
  citecolor=black,
  urlcolor=blue
}

\lstset{
  basicstyle=\ttfamily\small,
  breaklines=true,
  frame=single,
  columns=fullflexible
}

% =========================
% ДАННЫЕ ДЛЯ ТИТУЛЬНОГО ЛИСТА
% =========================
\title{Эвристические и ML-подходы к решению задачи 3D Bin Packing problem}
\authorfull{Иртюга Илья Владимирович}
\author{Иртюга И. В.}

\faculty{фундаментальных наук}
\facultyshort{ФН}
\chair{математического моделирования}
\chairno{}
\group{ФН12-81Б}

\chief{Фетисов Д. А.}
\FirstConsultant{}
\inspector{}

\workyear{2026}

\setcontentsname{СОДЕРЖАНИЕ}
\setrefsname{СПИСОК ИСПОЛЬЗОВАННЫХ ИСТОЧНИКОВ}

\begin{document}
\ifusekprj
  \maketitle
\else
  \begin{titlepage}
    \makeatletter
    \centering
    {\large Курсовая работа\par}
    \vspace{1.5cm}
    {\bfseries\LARGE \@title \par}
    \vspace{1.5cm}
    {\large Автор: \@author\par}
    \vspace{0.5cm}
    {\large Группа: \@group\par}
    \vspace{0.5cm}
    {\large Кафедра: \@chair\par}
    \vspace{0.5cm}
    {\large Руководитель: \@chief\par}
    \vfill
    {\large \@workyear\par}
    \makeatother
  \end{titlepage}
\fi
\tableofcontents
\pagebreak

% =========================
\section{Введение}

Задача трехмерной упаковки прямоугольных объектов возникает в складской логистике, ритейле и транспортировке грузов. От качества укладки зависит число используемых паллет, плотность перевозки, устойчивость груза и риск повреждения товаров. В общем виде эта постановка относится к семейству задач bin packing, которые являются вычислительно сложными даже в более простых одномерных и двумерных вариантах~\cite{wiki_bin_packing,martello2000}. Практический интерес к задачам такого типа подтверждается как академическими обзорами методов упаковки~\cite{isa2016}, так и прикладными соревнованиями по трехмерной компоновке объектов~\cite{kaggle_sleigh}.

В данной работе рассматривается проект по решению прикладной задачи 3D bin packing для укладки товаров на паллеты, представленный в открытом GitHub-репозитории~\cite{project_repo}. В отличие от абстрактной геометрической задачи, здесь необходимо учитывать ограничения, близкие к реальной логистике: допустимые повороты, массу, высоту, пересечения, устойчивость размещения и хрупкость отдельных товаров.

В отчете рассматриваются три подхода к решению задачи:
\begin{itemize}
  \item \texttt{Greedy} --- эвристический метод;
  \item \texttt{LNS} --- метаэвристический метод;
  \item \texttt{GAN+GA} --- гибридный нейросетево-эволюционный метод.
\end{itemize}

\textbf{Цель работы} --- описать постановку задачи 3D bin packing для паллет, рассмотреть три выбранных подхода и сравнить их качество на имеющихся экспериментальных результатах.

\textbf{Для достижения цели решаются следующие задачи:}
\begin{enumerate}
  \item сформулировать входные данные, выход решения и основные ограничения;
  \item описать метрику качества, по которой сравниваются решения;
  \item составить план реализации исследования в рамках отчета;
  \item классифицировать методы \texttt{Greedy}, \texttt{LNS} и \texttt{GAN+GA};
  \item провести сравнение методов по итоговому score и компонентам метрики;
  \item сформулировать выводы о применимости рассмотренных подходов.
\end{enumerate}

\textbf{Объект исследования} --- алгоритмы построения трехмерной укладки коробок на паллету.

\textbf{Предмет исследования} --- эвристические, метаэвристические и гибридные методы решения задачи 3D bin packing, а также метрики оценки качества получаемых укладок.

\pagebreak

% =========================
\section{Постановка задачи}

\subsection{Входные данные}

Вход одной задачи состоит из описания паллеты и набора типов коробок. Для паллеты задаются длина, ширина, максимально допустимая высота укладки и максимальная масса груза. Для каждого типа коробки задаются размеры, масса одного экземпляра, количество экземпляров и дополнительные признаки, влияющие на допустимые размещения.

К важным признакам коробок относятся:
\begin{itemize}
  \item \texttt{strict\_upright} --- запрет на поворот, изменяющий исходную вертикальную ориентацию;
  \item \texttt{fragile} --- признак хрупкости, ограничивающий размещение тяжелых объектов сверху;
  \item \texttt{stackable} --- возможность или запрет ставить другие коробки поверх данного объекта.
\end{itemize}

Таким образом, входные данные задают не только геометрию объектов, но и часть физических требований к укладке.

\subsection{Выход решения}

Результатом работы алгоритма является список размещенных коробок. Для каждого экземпляра указываются координаты нижнего угла, фактические размеры после выбранного поворота и код ориентации. Если часть коробок разместить не удалось, они фиксируются отдельно как неразмещенные.

Основная цель алгоритма --- построить допустимую укладку с максимально высоким качеством. При этом нельзя рассматривать качество только как плотность заполнения: решение должно быть корректным с точки зрения габаритов, массы, устойчивости и ограничений хрупкости.

\subsection{Жесткие ограничения}

В рассматриваемой постановке решение считается допустимым, если выполняются следующие условия:
\begin{enumerate}
  \item каждая размещенная коробка полностью находится внутри габаритов паллеты;
  \item коробки не пересекаются по объему;
  \item суммарная масса груза не превышает допустимую массу паллеты;
  \item для коробок со строгой вертикальной ориентацией не используется запрещенный поворот;
  \item коробка, лежащая не на основании паллеты, имеет достаточную опору снизу;
  \item тяжелые объекты не размещаются непосредственно над хрупкими товарами.
\end{enumerate}

Условие опоры задает дополнительную физическую реалистичность. Если коробка находится выше пола паллеты, то значительная часть площади ее основания должна опираться на верхние грани нижележащих объектов. В проекте используется порог 60\%:
\[
\frac{S_{\text{support}}}{S_{\text{base}}} \ge 0.6.
\]

Это ограничение отсекает геометрически возможные, но практически неустойчивые размещения.

\pagebreak

% =========================
\section{Метрика качества}

Итоговая оценка решения задается взвешенной суммой четырех компонент:
\[
\text{score} =
0.50 \cdot U_{\text{vol}} +
0.30 \cdot C_{\text{items}} +
0.10 \cdot F_{\text{frag}} +
0.10 \cdot T_{\text{time}}.
\]

Здесь $U_{\text{vol}}$ обозначает утилизацию объема паллеты, $C_{\text{items}}$ --- долю размещенных экземпляров, $F_{\text{frag}}$ --- оценку соблюдения ограничений хрупкости, а $T_{\text{time}}$ --- временной балл.

Утилизация объема показывает, какая часть доступного объема паллеты занята размещенными коробками:
\[
U_{\text{vol}} = \frac{\sum_i V_i}{LWH},
\]
где $V_i$ --- объем размещенной коробки, а $L$, $W$ и $H$ --- габариты паллеты.

Покрытие заказа определяется как отношение числа размещенных экземпляров к общему числу экземпляров:
\[
C_{\text{items}} = \frac{N_{\text{placed}}}{N_{\text{total}}}.
\]

Оценка хрупкости уменьшается при нарушениях правила размещения тяжелых объектов над хрупкими:
\[
F_{\text{frag}} = \max(0,\; 1 - 0.05 \cdot n_{\text{viol}}),
\]
где $n_{\text{viol}}$ --- число таких нарушений.

Временной балл отражает практическую применимость решения. Быстрые алгоритмы получают больший вклад в итоговую оценку, а методы, требующие слишком большого времени, теряют часть score. Поэтому в данной задаче важно не только найти плотную укладку, но и сделать это за разумное время.

\pagebreak

% =========================
\section{План реализации исследования}

Для отчета выбран компактный план, рассчитанный на объем не более 20 страниц. Он ориентирован на ключевые элементы задачи и не включает подробное описание всех файлов проекта.

\begin{enumerate}
  \item Описать прикладную мотивацию задачи 3D bin packing для паллет.
  \item Формализовать вход, выход, ограничения и итоговую метрику качества.
  \item Кратко описать данные и экспериментальный контур, ссылаясь на репозиторий как на единый источник реализации.
  \item Рассмотреть три метода: \texttt{Greedy}, \texttt{LNS} и \texttt{GAN+GA}.
  \item Для каждого метода указать тип подхода, основную идею, сильные и слабые стороны.
  \item Сравнить три метода по итоговому score на трех наборах задач.
  \item Проанализировать компоненты метрики и объяснить, почему один метод показывает лучший результат.
  \item Сформулировать итоговые выводы и возможные направления дальнейшего развития.
\end{enumerate}

Из рассмотрения намеренно исключаются дополнительные методы и специальные режимы сравнения, поскольку цель текущей версии отчета --- ясно сопоставить три выбранных подхода и уложиться в заданный объем.

\pagebreak

% =========================
\section{Данные и экспериментальный контур}

Для экспериментов используются синтетические наборы задач, имитирующие заказы продуктового ритейла. В них присутствуют товары разных типов: тяжелые, хрупкие, ориентируемые, компактные и объемные. Такой состав позволяет проверять алгоритмы не только на плотность заполнения паллеты, но и на способность учитывать ограничения устойчивости и хрупкости.

В проекте рассматриваются несколько сценариев сложности: однородные тяжелые грузы, смешанные наборы, задачи с большим числом экземпляров, а также наборы с заметной долей хрупких товаров. Для сравнения методов используются три общих набора задач, которые далее обозначаются как набор 100, набор 200 и контрольный набор.

Экспериментальный контур проекта устроен следующим образом: каждый алгоритм строит решения для одинаковых входных задач, после чего для каждого решения рассчитываются итоговый score и отдельные компоненты метрики. В отчете используются усредненные значения по наборам задач. Такой подход снижает влияние единичных удачных или неудачных примеров и позволяет сравнивать методы как классы алгоритмов.

Основной источник сведений о реализации, данных и полученных результатах --- GitHub-репозиторий проекта~\cite{project_repo}. В тексте отчета ссылки на отдельные файлы репозитория не используются, чтобы сохранить единый стиль оформления источников.

\pagebreak

% =========================
\section{Рассматриваемые методы}

\subsection{Greedy}

\textbf{Тип подхода:} эвристический.

Жадный алгоритм строит укладку последовательно: коробки рассматриваются одна за другой, для каждой коробки перебираются допустимые ориентации и возможные позиции, после чего выбирается локально лучший вариант размещения. Обычно такой выбор стремится положить объект как можно ниже и ближе к началу координат. Идеологически метод близок к классическим эвристикам first fit decreasing и bottom-left-fill, применяемым в задачах упаковки~\cite{wiki_bin_packing}.

Главное достоинство \texttt{Greedy} --- скорость. Алгоритм быстро строит допустимое решение и хорошо подходит как baseline, то есть начальная точка сравнения для более сложных методов. Кроме того, его поведение сравнительно легко объяснить: каждое решение принимается на основе текущего состояния укладки.

Главный недостаток метода --- локальность. Если на раннем этапе алгоритм выбрал неудачное размещение, он не умеет существенно перестраивать уже созданную структуру. Поэтому на сложных наборах задач жадный подход часто уступает методам, которые способны возвращаться к ранее принятым решениям.

\subsection{LNS}

\textbf{Тип подхода:} метаэвристический.

Метод \texttt{LNS} (\textit{Large Neighborhood Search}) использует идею крупноокрестностного поиска. Сначала строится начальное решение, а затем алгоритм многократно разрушает и восстанавливает его части. На этапе разрушения из укладки удаляется часть объектов, а на этапе восстановления освободившееся пространство заполняется заново. Такой подход позволяет исправлять локально неудачные фрагменты без полного построения решения с нуля.

Идеи локального и крупноокрестностного поиска широко применяются для сложных комбинаторных задач, в том числе для трехмерной упаковки~\cite{faroe2003,martello2000}. В данной постановке LNS особенно полезен, потому что качество укладки часто зависит от ранних решений: если метод может пересмотреть часть структуры, он получает шанс повысить утилизацию объема и покрытие заказа.

Преимущество \texttt{LNS} --- хороший баланс между качеством и временем. Он обычно работает дольше жадного алгоритма, но за это получает более плотные и полные укладки. Недостаток метода состоит в большей сложности настройки: результат зависит от того, какие фрагменты решения разрушаются, как выполняется восстановление и какой вычислительный бюджет доступен.

\subsection{GAN+GA}

\textbf{Тип подхода:} гибридный нейросетево-эволюционный.

Метод \texttt{GAN+GA} объединяет генетический алгоритм и генеративную модель. Генетическая часть отвечает за поиск в пространстве перестановок коробок и их ориентаций: создается популяция решений, затем применяются операции отбора, скрещивания и мутации. Генеративная компонента используется для получения новых перспективных конфигураций и поддержания разнообразия поиска.

Подходы, сочетающие генеративные модели и эволюционные алгоритмы, рассматриваются в современной литературе как один из способов усилить поиск в сложных задачах упаковки~\cite{ganga_paper}. Их идея состоит в том, чтобы не только случайно изменять решения, но и постепенно учиться порождать более удачные варианты.

Преимущество \texttt{GAN+GA} --- потенциально более широкий поиск по пространству решений. Метод не ограничен одним жадным порядком размещения и может исследовать разные конфигурации. Недостатки --- повышенная сложность настройки, зависимость от параметров стохастического поиска и отсутствие гарантии, что более сложная модель стабильно превзойдет хорошо настроенную метаэвристику.

\pagebreak

% =========================
\section{Экспериментальное сравнение}

\subsection{Сравнение итогового score}

В таблице~\ref{tab:main-score} приведены средние значения итогового score для трех рассматриваемых методов на трех общих наборах задач.

\begin{table}[H]
\centering
\caption{Сравнение итогового score}
\label{tab:main-score}
\begin{tabular}{lccc}
\toprule
\textbf{Метод} & \textbf{Набор 100} & \textbf{Набор 200} & \textbf{Контрольный набор} \\
\midrule
Greedy & 0.6842 & 0.7120 & 0.6549 \\
GAN+GA & 0.7012 & 0.7327 & 0.7161 \\
LNS & \textbf{0.7111} & \textbf{0.7434} & \textbf{0.7298} \\
\bottomrule
\end{tabular}
\end{table}

Во всех трех случаях лучший итоговый score показывает \texttt{LNS}. Метод \texttt{GAN+GA} стабильно превосходит \texttt{Greedy}, но остается ниже \texttt{LNS}. Наиболее заметный разрыв наблюдается на контрольном наборе: LNS превосходит Greedy примерно на 0.075 пункта, а GAN+GA --- примерно на 0.014 пункта.

Это означает, что в текущих экспериментах крупноокрестностный поиск оказался наиболее практичным компромиссом. Он использует больше времени, чем жадный метод, но получает более качественную структуру укладки.

\subsection{Сравнение компонент метрики}

Для более детальной интерпретации результатов рассмотрим отдельные компоненты итоговой метрики. Значения приведены в таблице~\ref{tab:main-components}.

\begin{longtable}{p{0.17\textwidth}p{0.20\textwidth}cccc}
\caption{Компоненты метрики для трех методов}\label{tab:main-components}\\
\toprule
\textbf{Метод} & \textbf{Набор} & \textbf{Vol.} & \textbf{Coverage} & \textbf{Fragility} & \textbf{Time} \\
\midrule
\endfirsthead
\toprule
\textbf{Метод} & \textbf{Набор} & \textbf{Vol.} & \textbf{Coverage} & \textbf{Fragility} & \textbf{Time} \\
\midrule
\endhead
Greedy & Набор 100 & 0.5961 & 0.6204 & 1.0000 & 1.0000 \\
GAN+GA & Набор 100 & 0.6278 & 0.6271 & 1.0000 & 0.9910 \\
LNS & Набор 100 & 0.6384 & 0.6580 & 0.9600 & 0.9850 \\
Greedy & Набор 200 & 0.6228 & 0.6686 & 1.0000 & 1.0000 \\
GAN+GA & Набор 200 & 0.6700 & 0.6591 & 1.0000 & 1.0000 \\
LNS & Набор 200 & 0.6712 & 0.7055 & 0.9747 & 0.9865 \\
Greedy & Контрольный & 0.6155 & 0.4906 & 1.0000 & 1.0000 \\
GAN+GA & Контрольный & 0.7156 & 0.5540 & 1.0000 & 0.9205 \\
LNS & Контрольный & 0.7172 & 0.6036 & 0.9833 & 0.9175 \\
\bottomrule
\end{longtable}

Из таблицы видно, что \texttt{Greedy} имеет наилучшее время работы и стабильно сохраняет высокий балл по хрупкости, но уступает по утилизации объема и покрытию заказа. \texttt{GAN+GA} улучшает утилизацию объема по сравнению с жадным методом и также хорошо соблюдает ограничения хрупкости. Однако его преимущество не всегда переходит в максимальный итоговый score.

\texttt{LNS} немного уступает по временной компоненте и иногда по хрупкости, но выигрывает по наиболее весомым частям метрики: утилизации объема и доле размещенных товаров. Поскольку эти две компоненты суммарно дают 80\% итогового score, именно они определяют лидерство LNS.

\subsection{Обсуждение результатов}

Результаты показывают, что усложнение метода само по себе не гарантирует лучшего качества. \texttt{GAN+GA} использует более сложную гибридную схему, однако в текущей реализации не превосходит \texttt{LNS}. Это можно объяснить тем, что для задач упаковки важна не только широта поиска, но и точное управление локальными перестройками укладки.

\texttt{Greedy} остается полезным как быстрый baseline и источник начальных решений. \texttt{LNS} выглядит наиболее сбалансированным методом: он использует структуру greedy-решения, но умеет пересматривать его фрагменты и улучшать ключевые компоненты метрики. \texttt{GAN+GA} представляет интерес как исследовательское направление, особенно если в дальнейшем усилить обучение генеративной части и подобрать параметры эволюционного поиска.

\pagebreak

% =========================
\section{Заключение}

В работе была рассмотрена прикладная задача 3D bin packing для укладки товаров на паллету. В отличие от упрощенной геометрической постановки, данная задача учитывает ограничения на массу, габариты, повороты, устойчивость и хрупкость товаров. Поэтому качество решения определяется не только плотностью упаковки, но и полнотой размещения, корректностью укладки и временем вычисления.

Были рассмотрены три подхода. \texttt{Greedy} является быстрым эвристическим baseline и хорошо подходит для получения простого допустимого решения. \texttt{LNS} относится к метаэвристикам и улучшает начальную укладку за счет разрушения и восстановления ее частей. \texttt{GAN+GA} представляет собой гибридный нейросетево-эволюционный подход, использующий идеи генетического поиска и генеративного моделирования.

Экспериментальное сравнение показало, что на имеющихся наборах задач лучший итоговый score демонстрирует \texttt{LNS}. Его преимущество связано прежде всего с лучшей утилизацией объема и более высоким покрытием заказа. \texttt{GAN+GA} показывает конкурентные результаты и превосходит \texttt{Greedy}, но пока не дает стабильного преимущества над LNS. Жадный алгоритм остается важным благодаря скорости и простоте интерпретации.

Таким образом, для рассматриваемой постановки наиболее рациональным основным методом является \texttt{LNS}, а \texttt{Greedy} и \texttt{GAN+GA} выступают соответственно как базовая эвристика и перспективная исследовательская альтернатива. Текущая структура отчета рассчитана на компактное изложение и должна укладываться в объем до 20 страниц.

\pagebreak

% =========================
\makeatletter
\renewcommand{\refname}{\@refsname}
\makeatother
\begin{thebibliography}{99}

\bibitem{wiki_bin_packing}
Bin packing problem // Wikipedia: The Free Encyclopedia [Электронный ресурс]. URL: \url{https://en.wikipedia.org/wiki/Bin_packing_problem} (дата обращения: 27.04.2026).

\bibitem{isa2016}
Статья из журнала Института системного анализа РАН [Электронный ресурс]. URL: \url{http://www.isa.ru/aidt/2016-03/31_43.pdf} (дата обращения: 27.04.2026).

\bibitem{kaggle_sleigh}
Packing Santa's Sleigh // Kaggle [Электронный ресурс]: соревнование по задаче трехмерной упаковки. URL: \url{https://www.kaggle.com/competitions/packing-santas-sleigh/overview} (дата обращения: 27.04.2026).

\bibitem{project_repo}
Lanity6. \textit{AI-CHAMP-CU-x-ITMO-X5-Tech} [Электронный ресурс]: GitHub-репозиторий проекта. URL: \url{https://github.com/Lanity6/AI-CHAMP-CU-x-ITMO-X5-Tech} (дата обращения: 27.04.2026).

\bibitem{martello2000}
Martello S., Pisinger D., Vigo D. The Three-Dimensional Bin Packing Problem // \textit{Operations Research}. 2000. Vol. 48. No. 2. P. 256--267. DOI: \url{https://doi.org/10.1287/opre.48.2.256.12386}.

\bibitem{faroe2003}
Faroe O., Pisinger D., Zachariasen M. Guided Local Search for the Three-Dimensional Bin-Packing Problem // \textit{INFORMS Journal on Computing}. 2003. Vol. 15. No. 3. P. 267--283.

\bibitem{ganga_paper}
Zhang Y., Wang X., Liu Z. A Novel Approach for Solving 3D Bin Packing Problem by Integrating a Generative Adversarial Network with a Genetic Algorithm // \textit{Scientific Reports}. 2024. DOI: \url{https://doi.org/10.1038/s41598-024-56699-7}.

\end{thebibliography}

\end{document}
